\scanpage{12}
\mspage{11}

III.5

the fact that $M\otimes_A\operatorname{Hom}_A(M,A)$ is projective as a left $W^e$-module.
Since $\mu$ is a left $W^e$-isomorphism, we see that $W$ is projective as a
left $W^e$ module, i.e., $W$ is separable over $A$.

THEOREM III.4. If $\mathcal M$ is a faithful locally projective sheaf of
modules of finite type over $\mathcal A$, then $\underline{\operatorname{Hom}}_{\mathcal A}(\mathcal M,\mathcal M)$ is a locally separable
$\mathcal A$-algebra.

\uline{Proof:} We know from Theorem II.2 that $\underline{\operatorname{Hom}}_{\mathcal A}(\mathcal M,\mathcal M)=\Omega$
is a locally projective module of finite type. Since $\mathcal M$ has finite
presentation as an $\mathcal A$-module, $\Omega_x=\operatorname{Hom}_{\mathcal A_x}(\mathcal M_x,\mathcal M_x)$ for every $x$ in
$X$. Since $\mathcal M_x$ is a finitely generated faithful $\mathcal A_x$-module for every $x$,
$\Omega_x$ is a central separable $\mathcal A_x$-algebra for every $x$. Thus by
Theorem III.2, $\Omega$ is a locally separable $\mathcal A$-algebra.

THEOREM III.5. If $\mathcal L_1$ and $\mathcal L_2$ are locally separable over $\mathcal A$,
then $\mathcal L_1\otimes_{\mathcal A}\mathcal L_2$ is locally separable.

\uline{Proof:} We know that for every $x$ in $X$, $(\mathcal L_1\otimes_{\mathcal A}\mathcal L_2)_x\approx\mathcal L_{1_x}\otimes_{\mathcal A_x}\mathcal L_{2_x}$.
Since $\mathcal L_{1_x}$ and $\mathcal L_{2_x}$ are central separable over $\mathcal A_x$, we have, from
[4, Proposition 1.5]the fact that $\mathcal L_{1_x}\otimes_{\mathcal A_x}\mathcal L_{2_x}$ is central separable over
$\mathcal A_x$ for every $x$ in $X$. Thus Theorem III.2 tells us that $\mathcal L_1\otimes_{\mathcal A}\mathcal L_2$ is a
locally separable $\mathcal A$-algebra.
